\documentclass[12pt]{article}
\usepackage{geometry,fancyhdr,xr,hyperref,ifpdf,amsmath,rcs,shortcuts}
\usepackage{lastpage,longtable,color,paunits,amsmath}
\geometry{letterpaper,headsep=5mm,head=10mm,hmargin={20mm,30mm},bottom=10mm,top=5mm}


\pagestyle{fancy} 

\RCS $Revision: 1.5 $
\RCS $Date: 2002/01/09 03:50:54 $

\fancypagestyle{first}{
\lhead{A405}
\chead{Aerosol size distributions}
\rhead{page~\thepage/\pageref{LastPage}}
\lfoot{} 
\cfoot{} 
\rfoot{}
}

\ifpdf
    \usepackage[pdftex]{graphicx} 
    \usepackage{hyperref}
    \pdfcompresslevel=0
    \DeclareGraphicsExtensions{.pdf,.jpg,.mps,.png}
\else
    \usepackage{hyperref}
    \usepackage[dvips]{graphicx}
    \DeclareGraphicsRule{.eps.gz}{eps}{.eps.bb}{`gzip -d #1}
    \DeclareGraphicsExtensions{.eps,.eps.gz}
\fi
\externaldocument[lec6,]{lec6}

\begin{document}

\pagestyle{first}

\vspace{0.1in}
\noindent
\textbf{Aerosol size distributions}
\vspace{0.1in}

Lohmann et al. Chapter 5  give several examples of histograms of aerosol
number concentration as a function of aerosol radius (e.g.~Figs 5.2-5.6
5.10, 5.11) .  These curves are called aerosol number distributions,
and are written $n(r)$, with typical units of ($\ccm\, \mum^{-1}$)
(pronounced ``number per cubic centimeter per micron bin width''.
Note that the radius bins could be measured in cm, in which case the
units would look like \un{cm^{-4}}, but would still be pronounced
``number per cubic centimeter per centimeter bin width''.

The key to making sense of any histogram is to focus on the product of
the density times the bin width:

\begin{equation}
  \label{eq:prod}
  n(r)dr
\end{equation}
This has units of \ccm, and is the number of aerosols per \cc with radii
ranging from $r$ to $r+dr$.

Since aerosols span a radius range of several orders of magnitude, it is
usual to work with $\log r$ (I'll write $\log_{10} r=\log r$ and
$\log_e r = \ln r$ below).  We can multiply (\ref{eq:prod}) by $r$ to get

\begin{equation}
  \label{eq:prod}
  n(r)dr = r\,n(r) \frac{dr}{r} = r\,n(r) d\ln r
\end{equation}
Either way, this is still number of aerosols per \cc with radii
ranging from $r$ to $r+dr$ (or alternatively $\ln r$ to $\ln r + d\ln r$).

Most of the plots in W\&H represent this a little differently, using
the cumulative number distribution $N(r)$ (\ccm) defined as:

\begin{equation}
  \label{eq:cum}
  N(r) = \int_r^\infty n(r^\prime) dr^\prime
\end{equation}
which is read as ``the number concentration of aerosols with radii
larger than $r$''.   As we will see later in the chapter, this cumulautive size
distribution $N(r)$ is the natural way to report readings from a cloud condensation
nuclei counter.

$N(r)$ and $n(r)$ are related by the one-dimensional version of Leibniz' rule:

\begin{equation}
  \label{eq:lieb1}
  \frac{d}{dr} \int_r^\infty n(r^\prime) dr^\prime = -n(r)
\end{equation}

Equation (\ref{eq:lieb1}) is just a consequence of the fundamental theorem of calculus:

\begin{equation}
  \label{eq:fund}
   \frac{d}{dr}\left [ \int_a^b f(r^\prime) dr^\prime \right ] = \frac{d}{dr}\left [ F(b) - F(a) \right ]
\end{equation}
and

\begin{equation}
  \label{eq:fund2}
   \frac{dF(r)}{dr} =  f(r)
\end{equation}
along with the fact that $n(\infty)=0$ if there is a finite number of aerosols.


Given the above we can see that a plot of $-dN/d\log r$ vs $d\log r$ still
gives the same result as (\ref{eq:prod}):

\begin{eqnarray}
  \label{eq:same}
  -\frac{dN}{d\log r} \times d\log r &=& -\ln 10 \frac{dN}{d\ln r} \times 
      \frac{1}{\ln 10} d \ln r = - r \frac{dN}{d r} d \ln r \nonumber\\
        & = & r n(r) \frac{dr}{r} = n(r) dr
\end{eqnarray}

Given that  $dr$, the width of the bin, is infinitesimally thin, we can
consider the radius to be constant in the bin, and find various moments
for the size range $r \rightarrow r + dr$:

Aerosol surface area per \cc:

\begin{equation}
  \label{eq:surface}
dS=  4 \pi r^2 n(r) dr
\end{equation}

Aerosol volume per \cc:

\begin{equation}
  \label{eq:volume}
dV  = \frac{4}{3} \pi r^3 n(r) dr
\end{equation}

Aerosol mass per \cc:

\begin{equation}
  \label{eq:mass}
dM  = \rho_{aero} \frac{4}{3} \pi r^3 n(r) dr
\end{equation}
where $\rho_{aero}$ is the aerosol density.

If we divide (\ref{eq:surface}), (\ref{eq:volume}) or (\ref{eq:mass})
by the total number concentration $N(0)=N_0$ then we get the average
surface area, volume or mass of an individual aerosol given the size distribution
$n(r)$.  One particularly
useful moment is the \textit{effective radius} $r_{eff}$, defined as the average
radius weighted by the aerosol surface area:

\begin{equation}
  \label{eq:reff}
  r_{eff}= \frac{\int_0^\infty r \left [ 4\pi \left ( \frac{r}{2} \right )^2 n(r) \right ] dr}%
              {\int_0^\infty  \left [ 4\pi \left ( \frac{r}{2} \right )^2 n(r) \right ] dr}
       = \frac{\left < r^3 \right >}{\left < r^2 \right >}
\end{equation}
where:

\begin{equation}
  \label{eq:bracket}
\left < x \right > = \frac{\int_0^\infty  x n(r) dr}{\int_0^\infty  n(r) dr} =  \frac{\int_0^\infty  x n(r) dr}{N_{tot}}
\end{equation}
Because the total surface area per unit mass plays a large role in determining the reflectance
of the aerosol size distribution, $r_{eff}$ is used frequently in studies of
the radiative effect of aerosols on the planetary energy balance.
\end{document}

%%% Local Variables:
%%% mode: latex
%%% TeX-master: t
%%% End:
